% from ega3-7-fr.tex
\oldpage[III]{211}
\begin{proposition}[7.9.11]
\label{III.7.9.11.fr}
Les hypothèses étant celles de (7.9.4), soit $\mathscr L$ un
$\mathscr O_X$-Module inversible très ample pour $Y$, et posons
$\mathscr P_\bullet(n)=\mathscr P_\bullet\otimes_{\mathscr O_X}
\mathscr L^{\otimes n}$ pour tout $n\in\mathbf Z$. Alors, pour tout $y\in Y$,
la fonction
\[
  n\longmapsto\operatorname{EP}(f,\mathscr P_\bullet(n);y)
  \tag{7.9.11.1}\label{III.7.9.11.1.fr}
\]
est un polynôme à coefficients dans $\mathbf Q$, qui est le même pour tous les
points d'une même composante connexe de $Y$.
\end{proposition}

\begin{proof}
Il est clair que $\mathscr P_\bullet(n)$ est un complexe de
$\mathscr O_X$-Modules $Y$-plats. En vertu de (7.9.6.2), on peut se borner au cas
où $\mathscr P_\bullet$ est réduit à un seul terme $\mathscr F\ne0$ de degré 0 ;
en outre, comme il s'agit de questions locales sur $Y$, on peut supposer $Y$
affine et $f$ projectif (II, 5.5.3) ; posons $X_y=f^{-1}(y)$, et soit
$\mathscr L_y=\mathscr L\otimes_{\mathscr O_Y}\kappa(y)$, qui est un
$\mathscr O_{X_y}$-Module très ample (II, 4.4.10) ; en vertu de (7.7.2), on a,
pour le foncteur $\mathscr T_\bullet$ relatif au complexe
$\mathscr P_\bullet(n)$,
\[
  \mathscr T_i(\kappa(y))
  =H^{-i}(X_y,\mathscr F_y\otimes\mathscr L_y^{\otimes n})
  \qquad
  \bigl(\text{où }\mathscr F_y=\mathscr F\otimes_{\mathscr O_Y}\kappa(y)\bigr) ;
\]
d'où résulte que $\operatorname{EP}(f,\mathscr F(n);y)$ n'est autre que la
caractéristique d'Euler-Poincaré $\chi_{\kappa(y)}(\mathscr F_y(n))$ définie dans
(2.5.1) ; le fait que (7.9.11.1) soit un polynôme résulte alors de (2.5.3) ; en
outre, pour chaque $n$, sa valeur est constante dans une composante connexe de
$Y$ (7.9.4), ce qui achève la démonstration.
\end{proof}

\begin{env}[7.9.12]
\label{III.7.9.12.fr}
Nous désignerons par $\operatorname{PH}(f,\mathscr P_\bullet;y)$ ou
$\operatorname{PH}(\mathscr P_\bullet;y)$ le polynôme (7.9.11.1), à coefficients
rationnels, et nous dirons que c'est le \emph{polynôme de Hilbert en $y$
relatif à $\mathscr P_\bullet$, $f$ et $\mathscr L$} (ou simplement le
\emph{polynôme de Hilbert en $y$ de $\mathscr P_\bullet$}, ou de $f$, s'il n'en
résulte pas de confusion) ; lorsque $Y$ est connexe non vide, on supprime la
mention de $y$ dans la notation et la terminologie. L'invariant ainsi obtenu
jouera un rôle essentiel au chap. V, dans la théorie des « modules » des
faisceaux cohérents quotients d'un faisceau cohérent donné.

Avec les notations de (7.9.6) et (7.9.11), on a
\[
  \operatorname{PH}(\mathscr P_\bullet;y)
  =\operatorname{PH}(\mathscr P'_\bullet;y)
  +\operatorname{PH}(\mathscr P''_\bullet;y)
  \tag{7.9.12.1}\label{III.7.9.12.1.fr}
\]
et en particulier
\[
  \operatorname{PH}(\mathscr P_\bullet;y)
  =\sum_i(-1)^i\operatorname{PH}(\mathscr P_i;y) ;
  \tag{7.9.12.2}\label{III.7.9.12.2.fr}
\]
cela résulte trivialement de (7.9.6.1) et (7.9.6.2). De même, avec les notations
et hypothèses de (7.9.7), on a
\[
  \operatorname{PH}(\mathscr P'_\bullet;y')
  =\operatorname{PH}(\mathscr P_\bullet;g(y')).
  \tag{7.9.12.3}\label{III.7.9.12.3.fr}
\]
\end{env}

La formule (7.9.12.2) ramène l'étude des polynômes de Hilbert d'un complexe à
celle des polynômes de Hilbert d'un seul $\mathscr O_X$-Module $Y$-plat. Ces
derniers admettent une interprétation remarquable indépendante de considérations
homologiques :

\begin{corollary}[7.9.13]
\label{III.7.9.13.fr}
Soient $Y$ un préschéma noethérien, $f:X\to Y$ un morphisme propre,
$\mathscr L$ un $\mathscr O_X$-Module inversible très ample pour $Y$,
$\mathscr F$ un $\mathscr O_X$-Module cohérent et $Y$-plat. Il existe un entier
$n_0$ tel que pour $n\geq n_0$, $f_*(\mathscr F(n))$ soit un
$\mathscr O_Y$-Module localement libre, de rang en $y\in Y$ égal à
$\operatorname{PH}(f,\mathscr F;y)(n)$.
\end{corollary}

\begin{proof}
Comme le morphisme $f$ est projectif (II, 5.5.3), il existe $n_0$ tel que pour
$n\geq n_0$

